% =============================================================================
% Capitulo 3 - Requisitos y diseño del sistema
% =============================================================================
\chapter{Requisitos y diseño del sistema}
\label{cap:requisitos_diseno}

\section{Requisitos del sistema}
\label{sec:requisitos}

\subsection{Requisitos funcionales}

El sistema propuesto debe cubrir el flujo completo desde la adquisición RF hasta la decisión final del detector. Los requisitos funcionales principales son los siguientes:

\begin{description}
  \item[RF1] Capturar muestras IQ en la banda de \SI{2.4}{\giga\hertz} mediante un receptor SDR, almacenando las capturas junto con sus metadatos de frecuencia central, tasa de muestreo y duración.
  \item[RF2] Generar espectrogramas tiempo-frecuencia de forma reproducible a partir de las muestras IQ, manteniendo constantes los parámetros de segmentación, STFT, escala de potencia y mapa de color.
  \item[RF3] Construir un dataset de detección de objetos con imágenes y etiquetas en formato YOLO, donde cada caja delimite una ráfaga activa del espectrograma.
  \item[RF4] Entrenar un detector YOLOv8 capaz de localizar regiones de señal y clasificarlas como \emph{drone} o \emph{interference}.
  \item[RF5] Evaluar el detector con particiones por captura para evitar fuga de información entre entrenamiento, validación y prueba.
  \item[RF6] Obtener métricas de detección, incluyendo precisión, recall y matriz de confusión normalizada.
  \item[RF7] Mantener una segunda etapa de identificación de modelos de dron, basada en un clasificador aplicado únicamente sobre regiones o espectrogramas previamente filtrados por la etapa de detección.
\end{description}

\subsection{Requisitos no funcionales}

Los requisitos no funcionales están orientados a la reproducibilidad y a la coherencia física de la representación:

\begin{description}
  \item[RNF1] Las particiones de entrenamiento, validación y prueba deben realizarse por captura completa, no por imagen individual.
  \item[RNF2] El procesamiento de espectrogramas debe ser determinista: misma ventana temporal, mismo tamaño de FFT, mismo solape, misma normalización y mismo mapa de color.
  \item[RNF3] Las aumentaciones aplicadas durante el entrenamiento deben respetar el significado físico de los ejes tiempo y frecuencia. No se deben aplicar inversiones, rotaciones ni transformaciones geométricas que alteren la interpretación RF.
  \item[RNF4] El sistema debe registrar los parámetros de entrenamiento y los artefactos generados, incluyendo pesos del modelo, curvas de aprendizaje, predicciones de validación y matrices de confusión.
  \item[RNF5] La evaluación debe distinguir entre errores de clasificación activa, confusión con \emph{background} y falsos positivos, ya que estas categorías tienen distinta interpretación operativa.
\end{description}

\section{Diseño del sistema}
\label{sec:diseno}

\subsection{Arquitectura general}
\label{subsec:arquitectura_general}

La arquitectura final se organiza como un pipeline en dos etapas. La primera etapa transforma la señal IQ en espectrogramas y aplica un detector YOLOv8 para localizar ráfagas activas. Cada detección proporciona una caja delimitadora, una clase y una confianza. La segunda etapa, planteada como extensión, toma las regiones detectadas como dron y aplica un clasificador multiclase para identificar el modelo de UAV.

\begin{center}
\begin{tabular}{p{0.22\textwidth}p{0.66\textwidth}}
\toprule
\textbf{Bloque} & \textbf{Función} \\
\midrule
Adquisición RF & Captura de muestras IQ con USRP en \SI{2.4}{\giga\hertz}. \\
Preprocesado & Segmentación en ventanas de \SI{0.1}{\second} y generación de espectrogramas STFT. \\
Auto-etiquetado & Detección de regiones brillantes y filtrado por tamaño para generar cajas YOLO. \\
Stage 1 & YOLOv8 detecta y clasifica ráfagas como \emph{drone} o \emph{interference}. \\
Stage 2 & Clasificador ResNet18 identifica el modelo de dron en regiones detectadas. \\
\bottomrule
\end{tabular}
\end{center}

La decisión de usar detección de objetos en la primera etapa se debe a una limitación observada en el enfoque de clasificación global. Un clasificador que recibe la imagen completa puede aprovechar información contextual no deseada, como el fondo de la captura o diferencias entre escenarios. En cambio, un detector de objetos se entrena con regiones locales y debe aprender la morfología de las ráfagas, que es la característica físicamente relevante.

\subsection{Diseño del dataset de espectrogramas}
\label{subsec:diseno_dataset}

Cada captura SigMF se divide en segmentos de \SI{0.1}{\second}. Esta duración permite observar patrones de ráfagas y saltos de frecuencia sin mezclar excesiva variabilidad temporal. Cada segmento se transforma en una imagen PNG de \num{1024}x\num{576} píxeles usando un mapa de color fijo. El dataset v3 queda formado por \num{3150} espectrogramas, distribuidos en entrenamiento, validación y prueba por captura:

\begin{table}[htbp]
\centering
\caption{Distribución del dataset v3 por partición.}
\label{tab:dataset_v3}
\begin{tabular}{lcc}
\toprule
\textbf{Partición} & \textbf{Imágenes} & \textbf{Capturas} \\
\midrule
Entrenamiento & 2100 & 14 \\
Validación & 450 & 3 \\
Prueba & 600 & 4 \\
\bottomrule
\end{tabular}
\end{table}

Las capturas con dron activo corresponden al DJI Mini 3 a distancias de \SI{5}{\meter} y \SI{10}{\meter}, con distintos niveles de congestión WiFi. Las capturas sin dron incluyen escenarios con WiFi y Bluetooth, de forma que la clase \emph{interference} no representa simplemente ruido, sino señales reales presentes en la banda.

\subsection{Diseño del auto-etiquetado de ráfagas}
\label{subsec:diseno_autolabel}

La generación manual de cajas delimitadoras sobre miles de espectrogramas sería costosa. Por ello se diseña un procedimiento de etiquetado débil basado en umbrales de energía y filtros morfológicos. El objetivo no es etiquetar toda energía presente, sino aislar pequeñas ráfagas compatibles con señales FHSS y descartar estructuras grandes, como barras WiFi anchas.

Los parámetros principales del auto-etiquetado son:

\begin{table}[htbp]
\centering
\caption{Parámetros del auto-etiquetado de cajas delimitadoras.}
\label{tab:parametros_autolabel}
\begin{tabular}{lll}
\toprule
\textbf{Parámetro} & \textbf{Valor} & \textbf{Interpretación} \\
\midrule
Percentil de umbral & 99,8 & Selecciona el 0,2\% más brillante. \\
Área mínima & 8 px & Permite ráfagas FHSS muy pequeñas. \\
Área máxima & 0,005 de la imagen & Descarta estructuras WiFi anchas. \\
Kernel morfológico & 1 & Evita fusionar blobs independientes. \\
Suavizado & 0,4 & Reduce ruido sin borrar ráfagas. \\
\bottomrule
\end{tabular}
\end{table}

Cada región detectada hereda la clase de la captura de origen. Las capturas \texttt{mini3\_drone} generan cajas de clase \emph{drone}, mientras que las capturas \texttt{mini3\_nodrone} generan cajas de clase \emph{interference}. Esta aproximación es débil, pero coherente con el protocolo experimental: en las capturas con dron el Mini 3 es la fuente FHSS principal, mientras que en las capturas sin dron los blobs pequeños proceden de Bluetooth, WiFi residual o ruido.

\subsection{Diseño del detector YOLO}
\label{subsec:diseno_yolo}

El detector elegido para la primera etapa es YOLOv8n, la variante ligera de YOLOv8. El modelo recibe espectrogramas redimensionados a \num{640}x\num{640} mediante \emph{letterbox}, y devuelve cajas delimitadoras con una clase y una confianza asociada. Las clases activas son:

\begin{itemize}
  \item \textbf{drone}: ráfagas asociadas a la firma OcuSync del DJI Mini 3.
  \item \textbf{interference}: ráfagas o estructuras pequeñas procedentes de Bluetooth, WiFi residual o ruido.
\end{itemize}

Las aumentaciones se configuran de forma conservadora. No se aplican rotaciones, inversiones verticales/horizontales, \emph{shear} ni perspectiva, ya que los ejes del espectrograma tienen significado físico: el eje horizontal representa tiempo y el eje vertical representa frecuencia. Alterarlos artificialmente podría generar ejemplos físicamente incoherentes.

\subsection{Diseño de la segunda etapa de identificación}
\label{subsec:diseno_identificacion}

La identificación de modelos de dron se mantiene como segunda etapa del TFG. En lugar de aplicarse directamente sobre cualquier espectrograma, se plantea que reciba como entrada únicamente regiones detectadas como \emph{drone} por la primera etapa, o espectrogramas completos previamente filtrados por presencia de dron. De esta forma, la arquitectura responde a dos preguntas distintas:

\begin{enumerate}
  \item ¿Existe una señal compatible con dron en el espectrograma?
  \item Si existe, ¿a qué modelo de dron se parece su firma RF?
\end{enumerate}

Para esta etapa se mantiene el uso de clasificadores convolucionales como ResNet18 \cite{he2016resnet}, entrenados sobre espectrogramas o recortes asociados a distintos modelos de UAV. La distancia se conserva como variable complementaria, no como clase principal.

\subsection{Diseño de la evaluación}
\label{subsec:diseno_evaluacion}

La evaluación del Stage 1 se realiza con métricas propias de detección de objetos: precisión, recall, mAP@50 y mAP@50-95. Además, se analiza la matriz de confusión normalizada para distinguir entre tres tipos de comportamiento:

\begin{itemize}
  \item detecciones correctas de \emph{drone} o \emph{interference};
  \item confusión cruzada entre las dos clases activas;
  \item errores hacia o desde \emph{background}, que representan falsos negativos o falsos positivos.
\end{itemize}

En una aplicación de detección de drones, la interpretación de estos errores no es simétrica. Un falso positivo de dron genera una alerta innecesaria, mientras que un falso negativo puede implicar no detectar una amenaza. Por ello, además del valor numérico de las métricas, se analiza la dirección de los errores.
